\paragraph{平凡通道窗口、Schur--Plancherel 极限与 generic pressure fan 比值律}
在固定阶 Schur 包的严格逆变换、primitive 坐标的 Euler 分解以及分拆压力相图都已闭合之后，还可把最新一段 Schur 断层后段材料进一步压缩为如下六条结果：前两条刻画平凡通道窗口与 primitive 坐标之间的单位上三角整自同构；其后两条给出恒等扇区主导下的 Schur--Plancherel 极限及非平凡能量比例；最后两条把 generic pressure chamber 中的通道比较与 Perron-purity 下的 hook 湮灭压缩为有限 partition fan 上的显式角色比值与线性压力差。

\begin{theorem}[平凡通道与 primitive 坐标之间的单位上三角整自同构]\label{thm:xi-time-part63c-trivial-channel-unitriangular-primitive-automorphism}
\leanverified{paper\_xi\_time\_part63c\_trivial\_channel\_unitriangular\_primitive\_automorphism}
固定分辨率 \(m\)，并记
\[
T_q(m):=T_{q,(q)}(m)
\qquad(q\ge 1),
\]
其中 \(T_{q,(q)}(m)=h_q(D_m)\) 为平凡 Schur 通道。沿用命题
\ref{prop:pom-schur-trace-jacobi-trudi-euler-primitive-fiber}
中的 primitive 坐标
\[
b_n(m):=\frac{1}{n}\sum_{d\mid n}\mu(d)\,S_{n/d}(m),
\]
以及 Euler 分解
\[
H_m(z)=\prod_{n\ge 1}(1-z^n)^{-b_n(m)}=\sum_{q\ge 0}T_q(m)z^q.
\]
则对每个 \(q\ge 1\)，存在整系数多项式
\[
P_q\in \ZZ[b_1,\dots,b_{q-1}]
\]
使得
\[
\boxed{\;
T_q(m)=b_q(m)+P_q\bigl(b_1(m),\dots,b_{q-1}(m)\bigr).\;}
\]
因此，对任意 \(Q\ge 1\)，坐标变换
\[
(b_1,\dots,b_Q)\longmapsto (T_1,\dots,T_Q)
\]
是 \(\ZZ^Q\) 上的单位上三角多项式自同构。其逆变换可递归写为
\[
\boxed{\;
b_q(m)=T_q(m)-P_q\bigl(b_1(m),\dots,b_{q-1}(m)\bigr)
\qquad(1\le q\le Q).\;}
\]
\end{theorem}
\begin{proof}
由命题 \ref{prop:pom-schur-trace-trivial-channel-primitive-fiber-expansion}，
\[
T_q(m)=
\sum_{\substack{(k_n)_{n\ge 1}\subset\NN\\ \sum_{n\ge 1}n k_n=q}}
\prod_{n\ge 1}\binom{b_n(m)+k_n-1}{k_n}.
\]
在全部满足
\[
\sum_{n\ge 1}n k_n=q
\]
的多重指标中，唯一含有 \(b_q(m)\) 的可行项只能是
\[
k_q=1,
\qquad
k_n=0\quad(n\neq q),
\]
其贡献恰为
\[
\binom{b_q(m)}{1}=b_q(m).
\]
确实，若 \(k_q\ge 2\)，则 \(qk_q\ge 2q>q\) 不可行；若 \(k_q=1\) 且某个 \(k_n>0\)（\(n<q\)），则
\[
\sum_{r\ge 1}r k_r\ge q+n>q,
\]
同样不可能发生。故其余全部可行项都只涉及 \(\{b_1(m),\dots,b_{q-1}(m)\}\)。把这些项之和记为
\[
P_q\bigl(b_1(m),\dots,b_{q-1}(m)\bigr),
\]
便得到所示上三角展开。

由于对每个 \(q\) 都有
\[
T_q=b_q+\text{(仅依赖于 }b_1,\dots,b_{q-1}\text{ 的整系数多项式)},
\]
故该坐标变换在每一级上都具有对角元为 \(1\) 的上三角形式，可逐级回代唯一反解。于是它给出 \(\ZZ^Q\) 上的单位上三角多项式自同构。证毕。
\end{proof}

\begin{corollary}[有限平凡通道窗口对全部低阶 Schur 断层的完备充分性]\label{cor:xi-time-part63c-finite-trivial-window-complete-schur-statistic}
固定 \(Q\ge 1\)。则对任意分拆 \(\lambda\vdash q\le Q\)，存在整系数多项式
\[
\mathcal P_{q,\lambda}^{(Q)}\in \ZZ[T_1,\dots,T_Q]
\]
使得
\[
\boxed{\;
\frac{T_{q,\lambda}(m)}{f^\lambda}
=
\mathcal P_{q,\lambda}^{(Q)}\bigl(T_1(m),\dots,T_Q(m)\bigr).\;}
\]
反之，对每个 \(q\le Q\)，存在整系数多项式
\[
\mathcal R_q^{(Q)}\in \ZZ[T_1,\dots,T_q]
\]
使得
\[
\boxed{\;
b_q(m)=\mathcal R_q^{(Q)}\bigl(T_1(m),\dots,T_q(m)\bigr).\;}
\]
因此，对固定 \(Q\)，下列两组数据逐 \(m\) 严格等价：
\[
\{T_1(m),\dots,T_Q(m)\}
\quad\Longleftrightarrow\quad
\left\{\frac{T_{q,\lambda}(m)}{f^\lambda}:1\le q\le Q,\ \lambda\vdash q\right\}.
\]
\end{corollary}
\begin{proof}
由推论 \ref{cor:pom-schur-trace-all-channels-primitive-fiber-coordinate}，
\[
\frac{T_{q,\lambda}(m)}{f^\lambda}\in \ZZ[b_1(m),\dots,b_q(m)]
\qquad(\lambda\vdash q).
\]
再由定理 \ref{thm:xi-time-part63c-trivial-channel-unitriangular-primitive-automorphism}，每个 \(b_j(m)\) 都可递归写成 \(\ZZ[T_1,\dots,T_j]\) 中的多项式。将这些表达代入上式，即得
\[
\frac{T_{q,\lambda}(m)}{f^\lambda}
=
\mathcal P_{q,\lambda}^{(Q)}\bigl(T_1(m),\dots,T_Q(m)\bigr)
\]
的整系数多项式表达。

反向公式正是定理 \ref{thm:xi-time-part63c-trivial-channel-unitriangular-primitive-automorphism} 的回代构造；把其中右端的 \(b_1,\dots,b_{q-1}\) 继续用已恢复的 \(T_1,\dots,T_{q-1}\) 代换，即得 \(\mathcal R_q^{(Q)}\)。

于是，固定窗口 \(\{T_1,\dots,T_Q\}\) 与全部低阶 Schur 断层
\[
\left\{\frac{T_{q,\lambda}}{f^\lambda}:q\le Q\right\}
\]
互为整多项式重编码，故两组数据逐 \(m\) 严格等价。证毕。
\end{proof}

\begin{theorem}[恒等扇区严格主导时的 Schur--Plancherel 极限]\label{thm:xi-time-part63c-identity-sector-schur-plancherel-limit}
固定 \(q\ge 2\)，并假设恒等扇区严格主导：
\[
\frac{P(1)}{1}>\max_{2\le r\le q}\frac{P(r)}{r}.
\]
则存在 \(\varepsilon>0\)，使对任意 \(\lambda\vdash q\) 都有
\[
T_{q,\lambda}(m)
=
\frac{(f^\lambda)^2}{q!}\,\kappa_1^{\,q}\,e^{mqP(1)}
+O\!\left(e^{m(qP(1)-\varepsilon)}\right).
\]
因此，若定义归一化通道权
\[
\Pi_{m,q}(\lambda):=
\frac{T_{q,\lambda}(m)}{\sum_{\mu\vdash q}T_{q,\mu}(m)},
\]
则有极限
\[
\boxed{\;
\Pi_{m,q}(\lambda)\longrightarrow \frac{(f^\lambda)^2}{q!}
\qquad(m\to\infty).\;}
\]
\end{theorem}
\begin{proof}
由推论 \ref{cor:pom-identity-sector-collapse-all-schur-channels}，对每个 \(\lambda\vdash q\)，都有
\[
T_{q,\lambda}(m)
=
\frac{(f^\lambda)^2}{q!}\,\kappa_1^{\,q}\,e^{mqP(1)}
+O\!\left(e^{m(qP(1)-\varepsilon)}\right)
\]
对某个统一的 \(\varepsilon>0\) 成立。

对所有 \(\mu\vdash q\) 求和，并利用经典恒等式
\[
\sum_{\mu\vdash q}(f^\mu)^2=q!,
\]
得到
\[
\sum_{\mu\vdash q}T_{q,\mu}(m)
=
\kappa_1^{\,q}e^{mqP(1)}
+O\!\left(e^{m(qP(1)-\varepsilon)}\right).
\]
故
\[
\Pi_{m,q}(\lambda)
=
\frac{\frac{(f^\lambda)^2}{q!}\,\kappa_1^{\,q}e^{mqP(1)}+O(e^{m(qP(1)-\varepsilon)})}
{\kappa_1^{\,q}e^{mqP(1)}+O(e^{m(qP(1)-\varepsilon)})}.
\]
分子分母同除以 \(\kappa_1^{\,q}e^{mqP(1)}\) 并令 \(m\to\infty\)，即得
\[
\Pi_{m,q}(\lambda)\to \frac{(f^\lambda)^2}{q!}.
\]
证毕。
\end{proof}

\begin{corollary}[Schur--Plancherel 能量与非平凡方差的普适极限比例]\label{cor:xi-time-part63c-schur-plancherel-energy-variance-ratio}
沿用定理 \ref{thm:xi-time-part63c-identity-sector-schur-plancherel-limit} 的假设。定义
\[
E_q(m):=
\frac{1}{q!}\sum_{\sigma\in S_q}\mathcal C_\sigma(m)^2
=
\sum_{\lambda\vdash q}\left(\frac{T_{q,\lambda}(m)}{f^\lambda}\right)^2,
\]
\[
V_q(m):=
\frac{1}{q!}\sum_{\sigma\in S_q}\bigl(\mathcal C_\sigma(m)-\overline{\mathcal C}(m)\bigr)^2
=
\sum_{\lambda\vdash q,\ \lambda\neq(q)}\left(\frac{T_{q,\lambda}(m)}{f^\lambda}\right)^2,
\]
其中
\[
\overline{\mathcal C}(m):=\frac{1}{q!}\sum_{\sigma\in S_q}\mathcal C_\sigma(m)=T_{q,(q)}(m).
\]
则有精确渐近
\[
\boxed{\;
E_q(m)\sim \frac{\kappa_1^{2q}}{q!}\,e^{2mqP(1)},
\qquad
V_q(m)\sim \frac{q!-1}{q!^2}\,\kappa_1^{2q}e^{2mqP(1)}.\;}
\]
特别地，
\[
\boxed{\;
\frac{V_q(m)}{E_q(m)}\longrightarrow 1-\frac{1}{q!}.\;}
\]
\end{corollary}
\begin{proof}
由定理 \ref{thm:pom-schur-plancherel-energy-identity} 与定理 \ref{thm:pom-schur-variance-decomposition}，\(E_q(m)\) 与 \(V_q(m)\) 分别由全部 Schur 通道与全部非平凡 Schur 通道的平方和给出。

再由定理 \ref{thm:xi-time-part63c-identity-sector-schur-plancherel-limit}，
\[
\frac{T_{q,\lambda}(m)}{f^\lambda}
=
\frac{f^\lambda}{q!}\,\kappa_1^{\,q}e^{mqP(1)}
+O\!\left(e^{m(qP(1)-\varepsilon)}\right).
\]
平方后得到
\[
\left(\frac{T_{q,\lambda}(m)}{f^\lambda}\right)^2
=
\frac{(f^\lambda)^2}{q!^2}\,\kappa_1^{2q}e^{2mqP(1)}
+O\!\left(e^{m(2qP(1)-\varepsilon)}\right).
\]
对 \(\lambda\vdash q\) 求和，并利用
\[
\sum_{\lambda\vdash q}(f^\lambda)^2=q!,
\]
得
\[
E_q(m)\sim \frac{\kappa_1^{2q}}{q!}\,e^{2mqP(1)}.
\]

对 \(V_q(m)\) 只需在求和中去掉平凡分拆 \((q)\)。由于
\[
f^{(q)}=1,
\qquad
\sum_{\lambda\vdash q,\ \lambda\neq(q)}(f^\lambda)^2=q!-1,
\]
故
\[
V_q(m)\sim \frac{q!-1}{q!^2}\,\kappa_1^{2q}e^{2mqP(1)}.
\]
最后相除即得
\[
\frac{V_q(m)}{E_q(m)}\to 1-\frac{1}{q!}.
\]
证毕。
\end{proof}

\begin{theorem}[generic pressure chamber 内任意两条 Schur 通道的严格指数比值律]\label{thm:xi-time-part63c-generic-pressure-fan-ratio-law}
固定 \(q\ge 2\)，并把压力向量写为
\[
p=(P(1),\dots,P(q))\in\RR^q.
\]
假设
\[
p\notin \bigcup\mathcal H_q,
\]
其中 \(\mathcal H_q\) 为定理 \ref{thm:pom-schur-pressure-fan-diagram} 的共振超平面族。则对每个 \(\lambda\vdash q\)，存在唯一主导分拆
\[
\nu_\lambda(p)\in \mathrm{Supp}(\lambda)
\]
与某个 \(\varepsilon_\lambda(p)>0\)，使
\[
T_{q,\lambda}(m)
=
f^\lambda\frac{\chi^\lambda(\nu_\lambda(p))}{z_{\nu_\lambda(p)}}K(\nu_\lambda(p))\,e^{m\mathcal P(\nu_\lambda(p))}
+O\!\left(e^{m(\mathcal P(\nu_\lambda(p))-\varepsilon_\lambda(p))}\right).
\]
由此得到：
\begin{enumerate}
  \item 若两条通道 \(\lambda,\mu\vdash q\) 共享同一主导分拆
  \[
  \nu_\lambda(p)=\nu_\mu(p)=:\nu,
  \]
  则
  \[
  \boxed{\;
  \frac{T_{q,\lambda}(m)}{T_{q,\mu}(m)}
  \longrightarrow
  \frac{f^\lambda\chi^\lambda(\nu)}{f^\mu\chi^\mu(\nu)}.\;}
  \]
  \item 若
  \[
  \mathcal P(\nu_\lambda(p))>\mathcal P(\nu_\mu(p)),
  \]
  则
  \[
  \boxed{\;
  \frac{1}{m}\log\left|\frac{T_{q,\lambda}(m)}{T_{q,\mu}(m)}\right|
  \longrightarrow
  \mathcal P(\nu_\lambda(p))-\mathcal P(\nu_\mu(p)).\;}
  \]
\end{enumerate}
\end{theorem}
\begin{proof}
第一段渐近正是定理 \ref{thm:pom-schur-pressure-fan-diagram} 的内容。

若 \(\nu_\lambda(p)=\nu_\mu(p)=\nu\)，则两式主指数相同，且共有因子
\[
\frac{K(\nu)}{z_\nu}e^{m\mathcal P(\nu)}
\]
在比值中完全约去，从而
\[
\frac{T_{q,\lambda}(m)}{T_{q,\mu}(m)}
=
\frac{f^\lambda\chi^\lambda(\nu)+O(e^{-m\varepsilon})}
{f^\mu\chi^\mu(\nu)+O(e^{-m\varepsilon})}
\]
对某个 \(\varepsilon>0\) 成立。由于 \(\nu\in\mathrm{Supp}(\lambda)\cap\mathrm{Supp}(\mu)\)，主系数均非零，令 \(m\to\infty\) 即得第一条。

若
\[
\mathcal P(\nu_\lambda(p))>\mathcal P(\nu_\mu(p)),
\]
则
\[
T_{q,\lambda}(m)=C_\lambda e^{m\mathcal P(\nu_\lambda(p))}\bigl(1+O(e^{-m\varepsilon_\lambda})\bigr),
\]
\[
T_{q,\mu}(m)=C_\mu e^{m\mathcal P(\nu_\mu(p))}\bigl(1+O(e^{-m\varepsilon_\mu})\bigr),
\]
其中 \(C_\lambda,C_\mu\neq 0\)。因此
\[
\log\left|\frac{T_{q,\lambda}(m)}{T_{q,\mu}(m)}\right|
=
m\bigl(\mathcal P(\nu_\lambda(p))-\mathcal P(\nu_\mu(p))\bigr)+O(1),
\]
两边除以 \(m\) 并令 \(m\to\infty\) 即得第二条。证毕。
\end{proof}

\begin{theorem}[generic Perron-purity 强制唯一最大分拆为 \texorpdfstring{$(q)$}{(q)} 并导致全部非平凡 hook 通道湮灭]\label{thm:xi-time-part63c-generic-perron-purity-kills-all-hooks}
固定 \(q\ge 2\)，并作如下假设：
\begin{enumerate}
  \item \(A_q\) 原始；
  \item 其 Perron 根
  \[
  \rho_q=e^{P(q)}
  \]
  简单，且仅出现在平凡通道 \(\lambda=(q)\)，即对所有 \(\lambda\neq(q)\) 都有
  \[
  \rho(B_{q,\lambda})<\rho_q;
  \]
  \item 压力向量 \(p=(P(1),\dots,P(q))\) 满足
  \[
  p\notin \bigcup\mathcal H_q.
  \]
\end{enumerate}
则成立：
\begin{enumerate}
  \item 唯一全局最大化的分拆必为
  \[
  (q),
  \]
  等价地
  \[
  \boxed{\;
  P(q)>\mathcal P(\nu)\qquad(\forall\,\nu\vdash q,\ \nu\neq(q)).\;}
  \]
  \item 分拆共振间隙严格为正：
  \[
  \boxed{\;
  \Delta_{\mathrm{part}}(q)>0.\;}
  \]
  \item 对所有非平凡 hook 分拆
  \[
  \lambda=(q-k,1^k),
  \qquad
  1\le k\le q-1,
  \]
  都有
  \[
  \boxed{\;
  B_{q,\lambda}=0,
  \qquad
  T_{q,\lambda}(m)=0\ \ \forall m\ge 1.\;}
  \]
\end{enumerate}
\end{theorem}
\begin{proof}
设 \(\nu_*\) 为全局最大化分拆。若 \(\nu_*\neq(q)\)，则由 generic 假设
\[
p\notin\bigcup\mathcal H_q
\]
与定理 \ref{thm:pom-schur-pressure-fan-diagram} 可知，全体分拆压力 \(\{\mathcal P(\nu)\}_{\nu\vdash q}\) 两两不同，故
\[
\mathcal R_{\max}=\{\nu_*\}
\]
为单点。

现考虑共振层权函数
\[
w(\nu):=K(\nu)\mathbf 1_{\{\nu=\nu_*\}}.
\]
由于 \(\nu_*\neq(q)\)，该函数不是常值类函数。于是由命题 \ref{prop:pom-schur-resonant-layer-total-cancellation}，必存在某个非平凡分拆 \(\lambda\neq(q)\) 使对应 Fourier 系数
\[
b_\lambda=
\frac{\chi^\lambda(\nu_*)}{z_{\nu_*}}K(\nu_*)
\]
非零。特别地，\(\nu_*\in\mathrm{Supp}(\lambda)\)。

再由定理 \ref{thm:pom-schur-pressure-fan-diagram}，该 \(\lambda\) 通道满足
\[
T_{q,\lambda}(m)
=
f^\lambda\frac{\chi^\lambda(\nu_*)}{z_{\nu_*}}K(\nu_*)\,e^{m\mathcal P(\nu_*)}
+O\!\left(e^{m(\mathcal P(\nu_*)-\varepsilon_\lambda)}\right),
\]
故其主指数率为 \(\mathcal P(\nu_*)\)。由定理
\ref{thm:pom-schur-channel-radius-character-sieve}，
\[
\rho(B_{q,\lambda})=e^{\mathcal P(\nu_*)}.
\]
而 \(\nu_*\neq(q)\) 且全局最大，结合 generic 假设可知
\[
\mathcal P(\nu_*)>\mathcal P((q))=P(q),
\]
从而
\[
\rho(B_{q,\lambda})=e^{\mathcal P(\nu_*)}>e^{P(q)}=\rho_q.
\]
这与 Perron-purity 假设
\[
\rho(B_{q,\lambda})<\rho_q\qquad(\lambda\neq(q))
\]
矛盾。故唯一全局最大化分拆只能是 \((q)\)。这证明了第 \(1\) 条。

由于 generic 假设排除了与任何其他分拆的平手，故 \((q)\) 一旦成为唯一最大值，就自动得到
\[
\Delta_{\mathrm{part}}(q)=P(q)-\max_{\nu\neq(q)}\mathcal P(\nu)>0.
\]
第 \(2\) 条成立。

最后，第 \(3\) 条由定理 \ref{thm:pom-partition-gap-positive-kills-hooks} 立即推出：严格正的分拆间隙强制所有非平凡 hook 通道满足
\[
B_{q,\lambda}=0,
\]
等价地
\[
T_{q,\lambda}(m)=\Tr(B_{q,\lambda}^m)=0
\qquad(\forall\,m\ge 1).
\]
证毕。
\end{proof}

\noindent
综上，平凡通道窗口与 primitive 坐标之间并非仅有可逆关系，而是形成单位上三角的整系数多项式自同构；因此固定窗口内的全部 Schur 断层都不再携带超出平凡通道的新代数信息。进一步，在恒等扇区严格主导时，Schur 通道归一化后刚性收敛到对称群的 Plancherel 测度，而其非平凡总能量比例趋于纯表示论常数 \(1-\frac{1}{q!}\)。与此同时，在 generic pressure chamber 中，任意两条通道的全部渐近比较都压缩为主导分拆上的角色比值与线性压力差；一旦再叠加 Perron-purity，则唯一最大分拆被强制锁定为 \((q)\)，并由此导致全部非平凡 hook 通道整体湮灭。

\endinput
